% =============================================================================
% §1.7 重写草稿 v4（独立片段，未 patch 进主稿）
% 相对 v3，针对用户第二轮 7 条「内容缺失」级批评：
%   (1) 开篇：先抛科学问题（结构约束是否真助长期预测）+ 为何需实验，再交代安排；
%       不再以「本节给出实验设置」起句。
%   (2) 新增首小节「数据生成」（subsec:data-generation）：REMUS 100 仿真器、海流
%       外源、控制块切分、两阶段轨迹质量过滤、训练/验证划分与种子——全章原缺。
%   (3) 评估小节重构为「短期与长期」：短期=留出块/轨迹预测（block/heldout），
%       长期=自由递推（主）；二者考核不同能力。
%   (4) 评价指标去「内/外」歧义：改为「任务精度指标（数据空间）」与
%       「物理诊断量（模型空间）」，并点明评判标准 vs 机理佐证。
%   (5) 对比/消融模型加代号 M0–M5、A1–A4，表改为模型-主序、正文统一引用。
%   (6) 初值扰动去英文混杂（profile/floor/std/yaw/rpm → 档位/下限/标准差/偏航/转速），
%       并讲透：为何是初值扰动（滤波态残差）、哪三个设计维度、训练 vs 评估两角色。
% 事实核验：数据集 filter = filter_trajectories 两阶段质量过滤（物理有效性 +
%   激励百分位），非状态滤波，故「滤波后状态」只作初值扰动的部署语义、不挂到数据生成。
%   模型集 = 注册表 10 个：phnode_full + phnode_qforce/phnode_merged_force +
%   se3_momentum/se3_accel_blackbox + blackbox_fullstate + 4 个 ablate_*。
% 红线：三轴不混淆；噪声为初值/滤波态鲁棒性正则、非导航后验；不报告性能数值/排序/
%   「最重要结构」结论；不宣称 v4lite 优越；§1.7 不预锁 §1.8 证据子集。
% label sec:training-baselines-protocol 保留（§1.8 引用）；小节 label 主稿无外部引用，可自定。
% =============================================================================

\section{实验设置}
\label{sec:training-baselines-protocol}

第~\ref{sec:structured-continuous-model} 节的参数化在函数类层面保证了正定逆质量、半正定耗散与斜对称零功率耦合，第~\ref{sec:energy-power-properties} 节给出了静水机械核心的功率平衡；这些性质由构造保证、不随训练改变。然而，这些结构约束是否真的转化为更可靠的长期运动预测，并不能由构造本身回答：模型作为长期轨迹预测器的实际行为，取决于数值积分、状态转换与数据分布的共同作用，只能通过训练后的长期递推与诊断来检验。本章实验因此围绕一个问题展开——在受控仿真数据与给定模型族下，逐项结构约束对长期自由递推的稳定性与精度各起什么作用——并以一组结构完整度递减的受控对照来回答。本节规定回答该问题所需的实验设置：数据如何生成（第~\ref{subsec:data-generation} 小节）、训练目标（第~\ref{subsec:training-objective} 小节）、短期与长期评估方式（第~\ref{subsec:rollout-eval} 小节）、评价指标（第~\ref{subsec:metrics} 小节）、用于隔离各结构因素的对比模型与消融（第~\ref{subsec:models-ablations} 小节），以及作为鲁棒性检验的初值扰动（第~\ref{subsec:ic-perturbation} 小节）。各结构因素的定量作用与排序在第~\ref{sec:results-structure-evidence} 节报告，本节不预先给出任何性能结论。

\subsection{数据生成}
\label{subsec:data-generation}

实验数据由 REMUS~100 自主水下航行器的六自由度刚体仿真器生成。仿真器以体坐标动力学为核心，采用显式时间积分推进平移与角运动，姿态以四元数表示以避免欧拉角奇异；推进器转速与舵面命令、海流与深度上下文作为外源输入驱动运动。每条轨迹由随机生成的命令序列激励，并切分为若干控制块：块内控制命令与外源上下文保持分段常值，块与块之间更新，从而既覆盖丰富的受控机动，又与第~\ref{subsec:training-objective} 小节按控制块组织的训练目标一致。

海流以分段常值外源流场注入：每条轨迹随机采样流速与流向，流速上限取 \(0.5~\mathrm{m/s}\)；有海流样本依式~\eqref{eq:velocity-contract} 的速度契约维持总体速度与相对水速度的转换。本章主环境为有海流设置，并保留无海流设置作为对照。

生成的原始轨迹经两阶段质量过滤后入库：先剔除违反 \(\SO(3)\) 正交性或速度界限的物理无效轨迹，再按激励质量的百分位阈值保留信息量充足的样本，使数据集偏向能有效约束动力学的受激轨迹，而非低激励的近平凡运动。过滤后的有海流数据集含一千条轨迹、每条一百五十个控制块，控制块时长 \(0.2~\mathrm{s}\)、状态采样间隔 \(0.05~\mathrm{s}\)；按 \(0.8\) 的比例划分训练与验证集。数据生成与模型训练各自记录随机种子，跨模型比较固定同一数据镜像，使模型间差异不掺入数据采样差异。仿真器的水动力系数与执行器参数取自 REMUS~100 的公开数据，仅用于生成轨迹并初始化机械先验，不作为可从轨迹唯一辨识的真值（第~\ref{sec:modeling-prelim} 节）。数据规模等定量设置另汇总于表~\ref{tab:experimental-config}。

\subsection{训练目标}
\label{subsec:training-objective}

训练以控制块为基本单位（第~\ref{subsec:data-generation} 小节）：在单个控制块内，控制命令与外源上下文保持分段常值，模型从块初值积分生成预测轨迹，
\begin{equation}
  y_0=\mathcal{T}_{d\to m}(s_0),\qquad
  \hat y(t_i)=\Phi_\theta(t_i;t_0,y_0),\qquad
  \hat s_i=\mathcal{T}_{m\to d}(\hat y(t_i)),
  \label{eq:block-prediction}
\end{equation}
其中 \(\mathcal{T}_{d\to m}\) 与 \(\mathcal{T}_{m\to d}\) 为式~\eqref{eq:data-to-model-state} 与式~\eqref{eq:model-to-data-state} 的数据态--模型态转换，\(\Phi_\theta\) 为式~\eqref{eq:complete-augmented-vector-field} 向量场经神经常微分方程积分得到的流映射。目标轨迹同样转换到模型态，使有海流样本的速度监督作用于相对水速度 \(\nu_r\)（依式~\eqref{eq:velocity-contract} 的速度契约），位置、姿态与执行器状态在两种表示中保持同一数值。加权训练目标取
\begin{equation}
  \mathcal{L}
  =
  \sum_i
  \lambda_x\norm{\hat x_i-x_i}^2
  +\lambda_R d_{\SO(3)}(\hat R_i,R_i)^2
  +\lambda_\nu\norm{\hat\nu_{r,i}-\nu_{r,i}}^2
  +\lambda_a\norm{\hat u_{a,i}-u_{a,i}}^2
  +\lambda_{\SO}\mathcal{L}_{\SO(3)} .
  \label{eq:training-loss}
\end{equation}
姿态误差以旋转测地距离 \(d_{\SO(3)}(\hat R,R)\) 度量、而非旋转矩阵的欧氏差，使监督直接作用于 \(\SO(3)\) 流形上的几何偏离，并对脱离旋转群的非物理偏移保持敏感；\(\mathcal{L}_{\SO(3)}\) 为旋转矩阵的正交性与行列式正则项。速度项以相对水速度为监督对象，与第~\ref{sec:state-current-contract} 节的速度契约一致；评估时预测轨迹经 \(\mathcal{T}_{m\to d}\) 回到数据态，终端位置、姿态与总体速度等可观测量在数据空间解释（第~\ref{subsec:metrics} 小节）。连续时间积分、优化器与网络容量等实现细节统一列于表~\ref{tab:experimental-config}。

\subsection{评估方式：短期与长期}
\label{subsec:rollout-eval}

块级训练目标只反映局部受控初值问题的拟合质量，无法说明误差在自由递推中的累积；因此评估分短期与长期两个层次，二者考核不同的能力。

短期评估在留出的控制块与轨迹上度量模型对局部动力学的逼近：以留出样本的块初值启动，积分一个或少数几个控制块，比较预测与真值的状态偏差。短期评估反映模型是否学到了正确的瞬时状态演化，是长期能力的必要而非充分条件——块级精度高的模型仍可能在自由递推中发散。

长期评估为本章主要评估方式，采用控制块串接的自由递推：首块由给定初值启动，其后每块以上一块的预测末状态为初值，并按数据协议更新位置原点、海流速度、控制命令与深度上下文。所得轨迹不再由真实状态校正，单块误差沿时间累积、放大，姿态漂移、速度发散与数值失败等仅在长期显现的问题由此暴露。长期评估使用受控仿真生成的 AUV 轨迹，区分无海流与有海流设置；有海流数据依式~\eqref{eq:velocity-contract} 维持总体速度与相对水速度的转换，使位姿几何按总体速度推进、水动力配对作用于相对水速度。

评估时域取 10、30 与 60~\(\mathrm{s}\) 三层，以 60~\(\mathrm{s}\) 为本章主时域；较短时域用于刻画误差随时间的增长形态与提前失败的位置，避免单一终端时刻掩盖中途已发散的轨迹。未完成目标时域的轨迹，其终端误差不计入该时域的条件误差统计，可比性由完成率反映。对未完成的轨迹，进一步区分模型预测发散与真值侧或数值侧的局部失败，使后续分析能把结构导致的长期不稳定与偶发的数值失败分开——这一区分对判断几何结构是否支撑长期稳定尤为关键。

\subsection{评价指标}
\label{subsec:metrics}

评价指标分两类，分别回答“模型预测得准不准”和“结构组件是否按设计工作”。第一类是任务精度指标，在数据空间定义，刻画模型作为长期轨迹预测器的可观测性能，模型的优劣以这类指标判定；第二类是物理诊断量，在模型空间定义，用于核验各结构组件的行为是否与设计一致，但不替代任务精度。两类层级不同：任务精度是评判标准，物理诊断是机理佐证。

任务精度以 60~\(\mathrm{s}\) 终端位置误差的中位数与第 95 百分位（P95）并列度量，分别刻画典型误差与尾部风险。采用并列双指标而非单一标量，是因为典型精度与尾部稳健并不总是同向：一个中位数很低的模型仍可能在少数轨迹上累积显著更大的误差，而单一中位数或单一均值都会掩盖这一点；其必要性将在第~\ref{sec:results-structure-evidence} 节由中位数与尾部排序不一致的具体情形体现。完成目标时域的轨迹另报告深度误差、旋转测地误差与总体线速度、角速度误差。完成率不参与精度排序，仅用于判断终端误差是否具可比性——终端误差只有在完成目标时域的样本上才有意义。

物理诊断量核验结构组件是否按设计工作：相对水速度误差核验 \(\nu_r\) 是否与第~\ref{sec:state-current-contract} 节的速度契约一致；机械能量跨度与最大能量变化反映机械存储是否在递推中保持有限量级；\(\SO(3)\) 行列式与正交性误差反映姿态矩阵是否保持接近旋转群。各量先在单个运行内对评估轨迹取统计量，再在随机种子之间取均值，记号 \(N\) 表示参与聚合的有效种子数。

\subsection{对比模型与结构消融}
\label{subsec:models-ablations}

结构约束的作用通过一组受控对照检验。对比模型沿结构完整度递减排列（表~\ref{tab:model-comparison-roles}），使每一步移除恰好对应一个结构因素：完整结构化模型 AUVHamNODE（M0）同时具备 \(\SE(3)\) 位姿几何、机械存储与 \(D_\theta\)/\(J_\theta\)/\(\tau_\theta\) 分解；端口哈密顿风格基线（M1、M2）保留机械存储核心，而改变保守势能或非保守广义力的组织方式；\(\SE(3)\) 动量黑箱与加速度黑箱（M3、M4）保留位姿几何乃至质量映射，却把动量或加速度动力学交由黑箱；全状态黑箱（M5）则不含任何显式几何或能量结构。由 M0 到 M5 的链条因而把显式位姿几何、机械能量核心与非保守力分解等结构因素逐层剥离，使各因素的作用可在对照中分离。其中 M1--M4 参照 \(\SE(3)\) 哈密顿及李群端口哈密顿神经常微分方程的思路~\citep{duong2021se3hamnode,duong2024liegroupphnode}，但均在统一的 AUV 状态契约、数据协议与评估方式下重新实现为受控候选，而非复刻文献中的原始模型。为使对照反映结构差异而非容量差异，各候选采用容量匹配的网络宽度，海流任务统一以 REMUS~100 的质量矩阵与执行器先验初始化；族内差异因而可在容量可比的前提下归因于结构配置。

\begin{table}[htbp]
  \centering
  \caption{对比模型与结构消融及其检验目标}
  \label{tab:model-comparison-roles}
  \begin{tabularx}{\textwidth}{l>{\raggedright\arraybackslash}p{0.30\textwidth}>{\raggedright\arraybackslash}X}
    \toprule
    代号 & 模型 & 检验目标 \\
    \midrule
    M0 & 完整结构化模型 AUVHamNODE & 含速度契约、\(\SE(3)\) 运动学、机械存储、\(D_\theta\)/\(J_\theta\)/\(\tau_\theta\) 分解与执行器滞后的完整候选项。 \\
    M1 & 端口哈密顿位形广义力基线 & 保留机械存储核心，以一般位形广义力替代标量势能诱导的保守力，检验能量核心的作用。 \\
    M2 & 端口哈密顿合并广义力基线 & 保留机械存储核心，改变非保守广义力的组织方式（合并而非分解）。 \\
    M3 & \(\SE(3)\) 动量黑箱 & 保留位姿几何与常质量映射，动量动力学交由黑箱。 \\
    M4 & \(\SE(3)\) 加速度黑箱 & 保留位姿几何，加速度动力学交由黑箱、不保留显式质量映射。 \\
    M5 & 全状态黑箱 & 不含显式位姿几何与机械能量结构的非结构化状态导数模型。 \\
    \midrule
    A1 & 无质量先验消融 & 移除物理质量矩阵的先验初始化，保留正定质量参数化。 \\
    A2 & 对角阻尼消融 & 半正定耦合耗散矩阵 \(D_\theta(\xi)\) 限制为非负对角形式。 \\
    A3 & 无零功率（升力）耦合消融 & 移除斜对称零功率分支 \(J_\theta(\xi)\nu_r\)。 \\
    A4 & 广义力条件化消融 & 广义力分支 \(\tau_\theta\) 收缩为主要由执行器状态或命令驱动。 \\
    \bottomrule
  \end{tabularx}
\end{table}

结构消融（A1--A4）在完整模型 M0 上逐项移除或弱化单个结构组件，每个消融对应一条可检验的结构假设，并与第~\ref{sec:energy-power-properties} 节的相应性质相对接。无质量先验消融（A1）移除基于物理质量矩阵的先验初始化（同时保留正定质量参数化），检验质量先验对动量--速度配对学习与长期递推的作用；对角阻尼消融（A2）把半正定耦合耗散矩阵 \(D_\theta(\xi)\) 限制为非负对角形式，检验跨通道耗散耦合相对独立通道耗散的作用；无零功率耦合消融（A3）移除斜对称分支 \(J_\theta(\xi)\nu_r\)，检验升力或横向水动力类零功率耦合在候选模型族内的作用，但不等同于断言真实水动力升力项被唯一辨识；广义力条件化消融（A4）把广义力分支 \(\tau_\theta\) 收缩为主要由执行器状态或命令驱动、弱化相对速度与海流条件化，检验开放广义力通道条件化范围的作用。

所有候选在统一的数据、训练与数值实现配置下比较\footnote{连续时间积分基于 torchdiffeq~\citep{chen2018neuralode} 的固定步长四阶 Runge--Kutta 实现。}，主要设置汇总于表~\ref{tab:experimental-config}，命令级参数与运行元数据保留在各运行记录中。进入定量比较的运行须共享同一数据协议、代码环境与评估方式；出现训练崩溃或长期递推发散的运行移出精度聚合并单独披露，其中全部种子均发散者不报告中位数、而记为长期稳定性失败，以与精度不佳相区别。\(N\) 为排除后参与聚合的有效种子数；异常的判定与披露规则见第~\ref{sec:results-structure-evidence} 节。

\begin{table}[htbp]
  \centering
  \caption{长期递推实验的主要配置}
  \label{tab:experimental-config}
  \begin{tabularx}{\textwidth}{>{\raggedright\arraybackslash}p{0.18\textwidth}>{\raggedright\arraybackslash}X}
    \toprule
    配置维度 & 设置 \\
    \midrule
    数据规模 & 海流数据集含一千条轨迹、每条一百五十个控制块；控制块时长 \(0.2~\mathrm{s}\)，状态采样间隔 \(0.05~\mathrm{s}\)，训练与验证按 \(0.8\) 划分。 \\
    速度契约 & 数据态保存总体速度，训练前按式~\eqref{eq:velocity-contract} 转换为模型态相对水速度。 \\
    数值积分 & 固定步长四阶 Runge--Kutta，单精度浮点状态；训练、块级评估与长期递推共用同一求解器。 \\
    网络容量 & 各候选采用容量匹配的网络宽度；海流任务以 REMUS 质量矩阵与执行器先验初始化。 \\
    优化 & AdamW，学习率经线性预热后余弦衰减，至多三百个 epoch；批量、学习率端点与正则权重记录于运行配置。 \\
    递推评估 & 重采样初值，覆盖伪随机二进制序列（PRBS）、线性调频（CHIRP）与 Ornstein--Uhlenbeck 过程（OU）三类控制场景，每场景 30 条评估轨迹；评估时域取 10、30 与 60~\(\mathrm{s}\)，主时域 60~\(\mathrm{s}\)。 \\
    \bottomrule
  \end{tabularx}
\end{table}

\subsection{初值扰动与鲁棒性评估}
\label{subsec:ic-perturbation}

前述评估默认初值精确已知。但在部署中，模型消费的初始状态来自导航系统对航位推算与惯性测量的滤波估计，必然带有误差；评估模型在带噪初值下的稳健性，因而是长期预测可用性的一部分。初值扰动正为此设计：它只对 \(t=0\) 的初始状态施加误差，模拟滤波后状态估计的残差，其后递推过程不再注入噪声。需强调，该接口刻画的是模型实际消费变量的误差半径，而非完整导航误差后验——它不复现导航滤波器的联合协方差结构，只约束初值各通道的误差预算。

初值扰动在训练与评估两个阶段引入，承担两个不同角色，分析中不应混为一谈：训练阶段，带噪初值作为鲁棒性正则，决定模型是否在带噪初值上学习；评估阶段，带噪初值作为鲁棒性测试，决定模型在何种扰动下受考核。同一种扰动既可作为训练干预、也可作为评估干预，其含义由所处阶段决定。

扰动的设计分三个维度。其一是受扰通道：噪声由干净数据态转为干净模型态后，在相对水速度 \(\nu_r\)、姿态 \(R\)、执行器反馈 \(u_a\) 与海流 \(v_c\) 四个通道上采样零均值扰动，再启动递推；如此控制的是模型实际消费的变量，并在海流设置下依式~\eqref{eq:velocity-contract} 保持 \(R\)、\(\nu_r\) 与 \(v_c\) 的语义一致，避免数据态扰动看似不大、而模型态输入已被过度污染的失配。其二是扰动强度：评估按强度递增分为名义与退化两档，另设一档在名义随机扰动上叠加固定航向偏置的有偏配置，用以考核对系统性航向误差的稳健性；各通道预算以 REMUS~100 的航位推算与惯性量级为参照，列于表~\ref{tab:noise-budget}。其三是采样结构：噪声可逐控制块独立采样，也可由同一轨迹的各控制块共享一次实现，二者在训练与评估两端均可实例化。

各通道预算的取值方式使部分语义相对训练数据分布定义。相对速度噪声按通道自然尺度取 \(\sigma_i=\max(\sigma_i^{\min},\,\alpha\,\hat\sigma_i)\)，其中 \(\hat\sigma_i\) 为训练集该通道的标准差、\(\sigma_i^{\min}\) 为下限、\(\alpha\) 为档位系数。由此，名义与退化档的部分强度是相对训练数据分布定义的，而非完全固定的绝对物理规格——同一数据集内的跨模型比较因而合理，跨数据集时实际强度则可能漂移，各运行遂记录最终落地的预算以备追溯。

\begin{table}[htbp]
  \centering
  \caption{初值扰动各通道的档位噪声预算（REMUS~100 航位推算/惯性量级参照）}
  \label{tab:noise-budget}
  \small
  \begin{tabularx}{\textwidth}{>{\raggedright\arraybackslash}X cccc}
    \toprule
    通道（参数） & 名义训练 & 名义评估 & 退化评估 & 航向偏置评估 \\
    \midrule
    相对线速度 \(\nu_r\)（\(\alpha\)）& 0.03 & 0.05 & 0.10 & 同名义 \\
    相对角速度 \(\nu_r\)（\(\alpha\)）& 0.03 & 0.05 & 0.10 & 同名义 \\
    姿态（偏航主导小角 / rad）& 0.0035 & 0.0050 & 0.0120 & 名义 \(+\) 固定 0.015 偏置 \\
    海流 \(v_{cx},v_{cy}\) / \(v_{cz}\)（m/s）& 0.008 / 0.004 & 0.012 / 0.006 & 0.030 / 0.015 & 同名义 \\
    执行器 \(\delta_r,\delta_s\)（rad）/ 转速（r/min）& 0.002 / 3 & 0.003 / 5 & 0.008 / 15 & 同名义 \\
    \bottomrule
  \end{tabularx}
  \par\vspace{2pt}
  \footnotesize 注：相对速度以 \(\sigma_i=\max(\sigma_i^{\min},\alpha\,\hat\sigma_i)\) 取值，\(\hat\sigma_i\) 为训练集标准差，表列 \(\alpha\)；线速度与角速度下限 \(\sigma_i^{\min}\) 分别为 \(0.005~\mathrm{m/s}\) 与 \(0.0015~\mathrm{rad/s}\)。航向偏置评估在名义评估的随机扰动上叠加固定 \(0.015~\mathrm{rad}\) 的偏航偏置，符号按样本固定、一条递推内不变。
\end{table}

训练端对噪声强度采用预热--爬坡--混合的渐进调度，以避免噪声初值在优化初期破坏训练稳定性，具体预热轮数、爬坡区间与混合比例记录于运行配置；评估端对所有模型固定噪声随机种子，使被扰动初值在不同模型之间可直接比较。
